\documentclass[12pt,a4paper]{article}

\usepackage[utf8]{inputenc}
\usepackage[T1]{fontenc}
\usepackage[ngerman]{babel}
\usepackage{amsmath,amssymb,amsthm,mathtools}
\usepackage[colorlinks=true,linkcolor=blue,citecolor=red,urlcolor=blue]{hyperref}
\renewcommand*{\HyperDestNameFilter}[1]{\jobname-#1}
\usepackage{geometry}
\geometry{margin=2.5cm}
\emergencystretch=1.5em
\usepackage{enumitem}
\usepackage{booktabs}
\usepackage{graphicx}
\usepackage{needspace}
\usepackage[capitalise,noabbrev,ngerman]{cleveref}

\newtheorem{theorem}{Theorem}[section]
\newtheorem{lemma}[theorem]{Lemma}
\newtheorem{proposition}[theorem]{Proposition}
\newtheorem{corollary}[theorem]{Korollar}
\newtheorem{conjecture}[theorem]{Vermutung}
\theoremstyle{definition}
\newtheorem{definition}[theorem]{Definition}
\theoremstyle{remark}
\newtheorem{remark}[theorem]{Bemerkung}

\title{Die Riemannsche Vermutung: Cross-Route-Synthese\\[4pt]
\large Teil~IV einer fünfteiligen Serie}
\author{Lukas Geiger\\
\small Bernau, Deutschland\\
\small ORCID: \href{https://orcid.org/0009-0005-7296-1534}{0009-0005-7296-1534}}
\date{\today}

\begin{document}
\overfullrule=0pt
\maketitle

\begin{abstract}\footnotesize
Dieses Paper, Teil~IV einer fünfteiligen Serie
\cite{Geiger2026partI,Geiger2026partII,Geiger2026partIII,Geiger2026partV},
sammelt drei Klassen von Material, die nicht in den Umfang von Teil~II
(Even Dominance) oder III (Spektrale Routen) passen: die Cross-Route-Synthese,
die alle drei aktiven Zweige~A, B, Z vergleicht; die ruhenden Routen und
externen Linien der weiteren Landschaft; sowie eine selbstkorrigierende
Dokumentation der Hypothesen-Korrekturen im iterativen Forschungsprozess.

Die Cross-Route-Synthese identifiziert ein gemeinsames Galerkin-Muster über
die Zweige hinweg --- jede bedingte Reduktion beinhaltet eine
Fixed-Parameter-Positivstruktur, die zu einer Uniform-Limit-Aussage über
einen wachsenden Parameter $N(\lambda)$, $d(\lambda)$ oder $m(\lambda)$
befördert werden muss --- und tabuliert die charakteristischen Konstanten
($C^* = 64\pi^2/3$ als aktueller Boundary-Moment-Benchmark für
$\lambda=5,7$ in Zweig~A, $\kappa_\infty \in [1{,}6,1{,}8]\times 10^{-3}$
für Zweig~B, $g_0(m) \approx 2\pi/\log\gamma_m$ für Zweig~Z).

\medskip\noindent
\textbf{Serien-Status (v3.0, 2026-05-26; Kritikalitäts-Linsen-Update 2026-05-31).}
Die Serie wurde von einer Trilogie zu einer fünfteiligen Serie umstrukturiert;
der Status der bedingten Reduktion aus v2.2 ist unverändert. Eine Arbeitsphase 2026
ergänzt \cref{sec:sharpened-wall}: eine \emph{heuristische} Kritikalitäts-Linse
($\mathrm{RH}\Leftrightarrow\Lambda=0$), die die Routen strategisch neu gewichtet,
eine schärfere Lokalisierung der gemeinsamen Wand (die analytische Fortsetzung über
$\Re s=1$), eine berechenbare Diagnostik aus dem \emph{klassischen}
Laguerre--P\'olya/Herglotz-Kriterium sowie vier ehrliche Schließungen
\emph{konkreter Konstruktionen} (keine Routen, mehrere davon Selbstkorrekturen
intern vorgeschlagener Konstruktionen). Keine davon schließt RH, widerlegt eine
Route oder ändert den Status der bedingten Reduktion.
\end{abstract}

{\footnotesize
\noindent\textbf{MSC 2020:} 11M26, 11M36, 47A75\\
\textbf{Schlüsselwörter:} Riemannsche Vermutung, Cross-Route-Synthese,
Galerkin-Muster, ruhende Routen, Hypothesen-Korrekturen,
Landschaftsübersicht
}

\clearpage
\tableofcontents
\clearpage


% =============================================================================
\section{Einleitung}
\label{sec:intro}
% =============================================================================

Dieses Paper, Teil~IV einer fünfteiligen Serie, dient als die
,,Vermischtes und Kombination''-Komponente des Landschaftsatlas:
\begin{description}[font=\bfseries,leftmargin=2em]
  \item[Teil~I \cite{Geiger2026partI}:] Grundlagen und Obstruktionen.
  \item[Teil~II \cite{Geiger2026partII}:] Even Dominance (Zweig~A).
  \item[Teil~III \cite{Geiger2026partIII}:] Spektrale Routen
    (Zweige~B und~Z).
  \item[Teil~IV (dieses Paper):] Cross-Route-Synthese, ruhende Routen,
    externe Linien, methodische Sackgassen, Hypothesen-Korrekturen.
  \item[Teil~V \cite{Geiger2026partV}:] Conclusio.
\end{description}

Die drei aktiven Zweige~A (Even Dominance), B (Hadamard Strip / I-1
Normal-Family) und Z (CCM Microcluster) werden ausführlich in den
Teilen~II und~III entwickelt. Dieses Paper identifiziert die strukturellen
Muster, die sie teilen, überblickt die weitere Landschaft ruhender und
externer Routen und liefert eine transparente Dokumentation der
Hypothesen-Korrekturen während der iterativen Entwicklung.


% =============================================================================
\section{Cross-Route-Synthese}
\label{sec:cross-route}
% =============================================================================

\subsection{Charakteristische Konstanten}

Jeder aktive Zweig produziert eine charakteristische Konstante oder ein
Skalierungsgesetz, dessen Identifikation die bedingte Reduktion verankert:

\begin{center}
\renewcommand{\arraystretch}{1.3}
\resizebox{\textwidth}{!}{%
\begin{tabular}{l|l|l|l}
\textbf{Zweig} & \textbf{Konstante} & \textbf{Skalierungsgesetz} & \textbf{Theoretische Identifikation} \\\hline
A (Even Dom.) & $C^* = 64\pi^2/3 \approx 210{,}55$ & $\mathrm{sLI} \sim C^*\,\lambda/(NL)$ für den $\lambda=5,7$-Benchmark; $\lambda=3$ separat und $\lambda=11$ offen & $16\int_0^\infty t^2/\sinh^2(t/2)\,dt$ (Rückwärtskandidat; exakter Collar-Kern bekannt) \\
B (Hadamard) & $\kappa_\infty \in [1{,}6,1{,}8]\times 10^{-3}$ & sättigt mit $\lambda$ & Finite-Part-Residualkrümmung des Foundation-Lock \\
Z (CCM) & $g_0(m) \approx 2\pi/\log\gamma_m$ & logarithmisch abnehmend & Montgomery--Odlyzko (klassisch) \\
\end{tabular}}
\end{center}

\begin{remark}[Identifikationsstatus]
Die Konstante $C^* = 64\pi^2/3$ für Zweig~A ist ein Identifikationskandidat
für den $\lambda=5,7$-Benchmark. Der Audit vom 27.~Mai zeigt, dass die
einfache Endpoint-Kern-Faktorisierung falsch ist; der exakte Collar-Kern
lautet $K(t)=e^t\cosh(t)/(2\sinh^2t)$, sodass eine rigorose Herleitung über
die Eigenmoden-Asymptotik laufen muss. Der $\lambda=3$-Sweep wird als
Niedrigband-Ausreißer geführt, während der einzelne Punkt
$\lambda=11,N=200$ ($C_{\mathrm{eff}}\approx168$) ein offener
Diskriminator bleibt
(Teil~II~\cite{Geiger2026partII}, Bemerkung zum Herleitungskandidaten). Die
Konstante $\kappa_\infty$ für Zweig~B ist empirisch im konservativen Bereich
$[1{,}6,1{,}8]\times 10^{-3}$ etabliert (Teil~III~\cite{Geiger2026partIII}). Die Lücke
$g_0(m)$ für Zweig~Z beruht auf klassischer Paarkorrelationsheuristik.
\end{remark}


\subsection{Das gemeinsame Galerkin-Muster}

Trotz ihrer unterschiedlichen analytischen Mechanismen teilen alle drei
Zweige ein gemeinsames strukturelles Muster:

\begin{enumerate}
  \item \textbf{Fixed-Parameter-Positivstruktur:} Bei jedem endlichen
    $(\lambda, N)$ (bzw.\ $(\lambda, m)$ für Zweig~Z) ist die relevante
    Größe berechenbar und hat das korrekte Vorzeichen bzw.\ die korrekte
    Schranke.
  \item \textbf{Uniform-Limit-Anforderung:} Die bedingte Reduktion
    erfordert die Beförderung der Fixed-Parameter-Aussage zu einer
    uniformen Aussage für $\lambda \to \infty$ mit einem Parameter
    ($N(\lambda)$, $d(\lambda)$, $m(\lambda)$), der mitwächst.
  \item \textbf{Transientes vs.\ konvergiertes Regime:} Der Aufwärmparameter
    ist zweigabhängig. Für den GF1/$C^*$-Zweig ist die aktuelle diagnostische
    Skala $\xi=NL/\lambda$, nicht $N/L$: Die $\lambda=5,7$-Zeilen betreten
    das Benchmark-Fenster nahe $\xi\approx 50$, während der Punkt
    $\lambda=11,N=200$ ein offener Präkonvergenz-Diskriminator bleibt. Andere
    Galerkin-Zweige führen weiterhin ihre eigenen Auflösungsquotienten.
\end{enumerate}

Dieses Muster ist konsistent mit einer \emph{Bernstein--Galerkin-Vermutung}:
die Galerkin-Trunkierung der Weil-Quadratform bei Dimension $N$ konvergiert
gegen die unendlichdimensionale Form in der relevanten Operatortopologie,
und die Konvergenz ist uniform in $\lambda$ nach einer
$\lambda$-abhängigen Aufwärmphase.

\begin{remark}[Status der Vermutung]
Das Bernstein--Galerkin-Muster ist eine empirische Beobachtung aus
5--7 Datenpunkten pro Zweig. Es wurde nicht als mathematische Aussage
formalisiert. Seine Rolle ist heuristisch: es liefert eine mechanistische
Erklärung dafür, warum alle drei Zweige ähnliche strukturelle Lücken
aufweisen (Fixed-Parameter $\checkmark$, Uniform-Limit offen).
\end{remark}


\subsection{Zwei isolierte offene Beweisslots}

Die Cross-Route-Analyse identifiziert zwei konkrete, isolierte offene
Beweisslots, deren Lösung einen oder mehrere Zweige schließen würde:

\begin{enumerate}
  \item \textbf{B-1: Uniforme Log-Krümmungskontrolle aus dem Foundation-Lock.}
    Falls $H_\lambda(\delta) \leq K\delta^2$ mit
    $K \leq 1{,}8\times 10^{-3}$ aus der Foundation-Lock-Struktur
    hergeleitet wird, dann liefert Theorem~B-1
    (Teil~III~\cite{Geiger2026partIII})
    $A_\lambda \leq 1{,}0063$, was Zweig~B schließt.
  \item \textbf{GF1: Eigenmoden-Asymptotik + Faktor-16-Identifikation.}
    Falls der exakte Collar-Kern und die $q_a(k)$-Eigenmoden-Asymptotik
    rigoros den Faktor~$16$ in $C^* = 16\cdot 4\pi^2/3$ erzeugen und der
    $\lambda=11$-Diskriminator nicht als separater Ast persistiert, dann liefert
    Theorem~GF1 (Teil~II~\cite{Geiger2026partII}) einen
    Randmomenten-Benchmark für das nicht-ausnahmsartige Zweig-A-Regime;
    der $\lambda=3$-Ast bleibt ein konservativer Guardrail.
\end{enumerate}

Diese beiden Slots sind \emph{unabhängig}: die Lösung des einen beeinflusst
den anderen nicht. Jeder ist ein wohlgestelltes analytisches Problem, das mit
Standardwerkzeugen zugänglich ist (komplexe Analysis für B-1; Fourier-Analysis
und Spektraltheorie für GF1).

\begin{remark}[Beide Slots sind Marge-artig --- ein strategisches Caveat]
\label{rem:slots-margin}
Beide Slots sind \emph{Marge-/Koerzivitäts}-Aussagen: eine uniforme Schranke, die
eine Lücke schließt. Sie bleiben wohlgestellte offene Probleme, und nichts im
Folgenden widerlegt sie. \Cref{sec:sharpened-wall} ergänzt eine \emph{heuristische}
Kritikalitäts-Linse: da RH (falls wahr) beim kritischen Wert $\Lambda=0$ ohne Puffer
sitzt (\cref{sec:lambda-lens}), verdienen exakt-strukturelle Mechanismen (Identität,
Symmetrie, Spektralrealisierung) \emph{zusätzliches} strategisches Gewicht neben
diesen Marge-artigen Slots. Dies ist eine Neugewichtung der Prioritäten, kein
Theorem und keine Widerlegung (siehe das Caveat in \cref{rem:lambda-caveat}).
\end{remark}


% =============================================================================
\section{Die geschärfte gemeinsame Wand: eine Kritikalitäts-Linse und eine Analytische-Fortsetzungs-Lokalisierung}
\label{sec:sharpened-wall}
% =============================================================================

Eine Arbeitsphase 2026 untersuchte alle Zweige erneut unter einer Frage ---
\emph{welche Art von Argument kann die Lücke schließen?} --- und lokalisierte die
gemeinsame Obstruktion schärfer. Alles in diesem Abschnitt ist eine
\textbf{strategische Linse und eine diagnostische Lokalisierung}, keine neue
Schließung von RH und keine Widerlegung einer Route. Mehrere Punkte sind ehrliche
Selbstkorrekturen intern vorgeschlagener Konstruktionen. Der Status der bedingten
Reduktion der Teile~II/III ist unverändert.

\subsection{Eine Kritikalitäts-Linse (heuristisch, kein Theorem)}
\label{sec:lambda-lens}

Die de~Bruijn--Newman-Konstante $\Lambda$ erfüllt $\Lambda\le 0\Leftrightarrow$ RH;
Rodgers und Tao (\emph{Ann.\ of Math.}\ 2020) bewiesen $\Lambda\ge 0$. Somit
$\mathrm{RH}\Leftrightarrow\Lambda=0$: falls RH gilt, sitzt sie \emph{exakt} beim
kritischen Wert, ohne Marge.

\begin{remark}[Kritikalitäts-Linse --- und ihr essenzielles Caveat]
\label{rem:lambda-caveat}
Dies lädt zu einer \emph{strategischen} Dichotomie von Beweismechanismen ein:
\emph{Marge-suchend} (eine uniforme Lücken-/Koerzivitäts-/Smallness-Schranke) versus
\emph{exakt-strukturell} (eine Identität, eine Symmetrie oder eine
Spektralrealisierung, die die Schlussfolgerung ohne Spielraum erzwingt). Die
heuristische Lesart ist, dass ein Mechanismus, der eine strikte uniforme Marge
erzeugt, etwas mit Puffer etablieren würde, was schwer mit einer kritischen Wahrheit
bei $\Lambda=0$ vereinbar ist.
\textbf{Caveat (essenziell, gegen Overclaiming).} Der Schritt ,,eine uniforme Marge
würde $\Lambda<0$ erzwingen`` ist eine \emph{Analogie}, kein Theorem. Die
zweigspezifischen Margen --- die Even-Dominance-Lücke $\lambda_1^+<\lambda_1^-$, die
Hadamard-Konstante $\kappa$, die CCM-Residual-Smallness --- sind \emph{nicht} als
eine de~Bruijn--Newman-Aussage gezeigt; $\Lambda$ und diese Größen sind
verschiedene Objekte. Folglich widerlegt die Linse die wohlgestellten Beweisslots
aus \cref{sec:cross-route} \emph{nicht} (\cref{rem:slots-margin}); sie legt nur nahe,
exakt-strukturelle Routen \emph{zusätzlich} zu gewichten. Es ist eine Neugewichtung
der Prioritäten, kein Ergebnis.
\end{remark}

Unter dieser Linse sortieren sich die Routen nach \emph{Typ}, nicht nach
Tragfähigkeit: die Marge-artigen Zweige~A (A7/AVG-Lücke), B ((ii-a)-Koerzivität), Z
(Residual-Smallness) behalten ihre exakten Kerne als Assets, während die
exakt-strukturellen Linien --- M (Bost--Connes\,$\leftrightarrow$\,Meyer, mit dem
KMS-Phasenübergang bei $\beta=1$ in der Rolle eingebauter Kritikalität) und die
globale adelische Trace-Route (Connes--Consani-Scaling-Site / $\mathbb{F}_1$, wo das
einzig fehlende Datum die Hodge-Index-Positivität einer ansonsten verfügbaren
Lorentz-Form ist) --- zusätzliches strategisches Gewicht erhalten. Die empirische
Signatur ist konsistent: die Even-Dominance-Marge wächst wie ein Potenzgesetz
$\sim\lambda^{0{,}72}$ (eine Kritikalitäts-Signatur, kein Puffer). Diese
Neugewichtung ist im Projekt-Entscheidungslog festgehalten; sie ändert den formalen
Status keiner Route.

\subsection{Die gemeinsame Wand: analytische Fortsetzung über $\Re s=1$}
\label{sec:common-wall-ac}

Unabhängig von der Linse teilen die korrekt typisierten Routen einen einzigen
konkreten Endpunkt. Drei Arbeitslinien --- der Möbius-Primzahlhub, die
Wiener--Hopf-Faktorisierung und die Sonin/Toeplitz-Konstruktion (Connes--Consani)
--- sind jeweils \emph{lokal/endlich} kontrollierbar, brechen aber an \emph{exakt}
derselben Stelle: dem globalen Limes, wo sich die lokalen Daten zu $\zeta$
zusammensetzen müssen. Die einheitliche Beschreibung ist die
Weyl--Titchmarsh-Funktion $M_\infty$ des unbeschränkten Limes-Operators; das
fehlende Datum ist uniforme Kontrolle dieses Limes aus lokalen Daten \emph{ohne}
$\zeta$ zu importieren.

\begin{remark}[Die Wand ist ein klassisches Faktum, hier scharf lokalisiert]
\label{rem:wall-folklore}
Ein endliches Euler-Produkt $\prod_{p\in P}(1-p^{-s})^{-1}$ hat nach der Substitution
$z=-\mathrm{i}(s-\tfrac12)$ einen Pol-Kamm in der oberen Halbebene bei
$z_{p,k}=2\pi k/\log p+\mathrm{i}/2$ (auf $\Im z=\tfrac12$), mit einem gemeinsamen Pol
bei $z=\mathrm{i}/2$ der Ordnung $|P|$; folglich ist seine logarithmische Ableitung
keine Herglotz-Funktion. Der Kamm wird nur durch $\zeta$'s analytische Fortsetzung
entfernt ($\zeta(0)$ endlich; der Pol bei $s=1$). Dies ist
\textbf{Euler-Produkt-Folklore} (das Produkt konvergiert nur für $\Re s>1$). Es ist
keine neue Obstruktion: es ist eine \emph{präzise Lokalisierung} der bekannten Wand
--- das fehlende Datum ist die analytische Fortsetzung über $\Re s=1$, d.\,h.\ ein
globales/Funktionalgleichungs-Phänomen, das endliche arithmetisch-plus-archimedische
Daten strukturell nicht liefern.
\end{remark}

\subsection{Eine berechenbare Diagnostik aus dem klassischen Laguerre--P\'olya/Herglotz-Kriterium}
\label{sec:diagnostic}

Sei $\Xi(z)=\xi(\tfrac12+\mathrm{i}z)$ und $M_\Xi=-\Xi'/\Xi$. Die Kette
\[
\mathrm{RH}\ \Longleftrightarrow\ \Xi\in\text{Laguerre--P\'olya}
\ \Longleftrightarrow\ M_\Xi\ \text{Herglotz}
\ \Longleftrightarrow\ \tfrac{(M_\Xi(z)-\overline{M_\Xi(w)})}{z-\bar w}\succeq 0
\]
ist \textbf{klassisch} (Hermite--Biehler; die Laguerre--P\'olya-Charakterisierung
reell-nullstelliger ganzer Funktionen; die Positivität $\Re\,\xi'/\xi>0$ für
$\Re s>\tfrac12$ ist das Voros-Kriterium). \emph{Wir beanspruchen keine Neuheit für
das Kriterium selbst.} Der Beitrag ist ein \emph{berechenbares, nicht-zirkuläres
Test-Protokoll}, das diese klassische Äquivalenz in ein diagnostisches Instrument
verwandelt: ein Euler-Anker, der $M_\Xi$ aus $\xi'/\xi$ in der konvergenten Region
$\Re w>1$ liefert (numerisch auf $\sim 10^{-31}$ verifiziert), eine Leiter von Gates
(Normalisierung; Tail-Tightness; eine Stieltjes-Identität; eine Off-Axis-Poisson-
Tomographie, die einzelne Residuen auflöst), und ein Detector-Lemma, durch das eine
Off-Line-Nullstelle eine negative $1\times1$-Pick-Minore knapp unterhalb des
gewanderten Pols erzwingt. Das Instrument ist empirisch validiert: es detektiert
Off-Line-Nullstellen und unterscheidet das echte Nullstellen-Spektrum von einer
glatten, dichte-angepassten Kontrolle.

\begin{remark}[Was die Diagnostik an einem echten Operator etablierte]
Eine unabhängige, nicht-zirkuläre Diskretisierung des Connes--Moscovici-Prolaten-
Operators (arXiv:2112.05500) wurde durch das volle Instrument geführt. Sie
reproduziert Riemanns Nullstellen-Zählfunktion $N(E)$ auf $\sim 3\%$ im UV --- die
\emph{archimedische Hälfte} (die glatte Dichte, führender Term
$\tfrac12\log(y/2\pi)$), eine unabhängige nicht-zirkuläre Bestätigung der
CCM-Dichte-Behauptung. Aber ihre Pole sitzen $O(1)$ \emph{neben} den Nullstellen,
sodass die arithmetische Feinstruktur (die \emph{Positionen}) nicht erfasst wird; und
die relative-Determinante-/Krein-Konstruktion entfernt den Euler-Pol-Kamm
\emph{nicht}, weil der prolate Hintergrund selbst Herglotz (kammfrei) ist, sodass
nichts zu kürzen bleibt. Beide Befunde bestätigen \cref{rem:wall-folklore} erneut:
Dichte ja, Positionen nein; die Wand ist die analytische Fortsetzung. Das genuine
positive Nebenprodukt --- dass der Determinanten-Kanal $-D'/D$
(Residuen~$=$~Multiplizitäten), nicht bloße Eigenwert-Konvergenz, das Objekt ist, das
konvergieren muss ($\det_{\mathrm{reg}}(A-z)\to C\cdot\Xi$) --- schärft die
CCM-Frontlinie und ist unter den unabhängigen Ergebnissen festgehalten
(\cref{sec:independent-results}).
\end{remark}

\subsection{Vier Schließungen konkreter Konstruktionen (keine Routen)}
\label{sec:four-closures}

Die Phase brachte vier genuine Negativ-Resultate hervor, jedes über ein
\emph{spezifisches konstruiertes Objekt bzw.\ eine Technik}. Keines schließt eine
ganze Route; jedes re-typisiert die Route oder benennt einen Nachfolger.

\begin{enumerate}
  \item \textbf{Osterwalder--Schrader-Reflexion, natürliche Form.} Der
    Reflexionsblock $Q_+(x)=\iint_0^\infty \overline{x(t)}\,K(t+s)\,x(s)$ ist
    indefinit, weil die Prim-Distributionen $\delta(t+s-\log p)$ im Weil-Kern die
    Halbachse $t+s=\log p$ treffen; die Toeplitz-Alternative $K(t-s)$ kollabiert zu
    globaler Weil-Positivität $=$ RH. Die OS-Route in ihrer natürlichsten Form
    (Inversions-Reflexion auf $\{|g|\ge 1\}$) ist geschlossen-negativ; eine andere
    Halbraum-Wahl wäre nötig.
  \item \textbf{Möbius-$\mathbb{Z}_2$-Primzahlhub, naives $L_s^-$.} Ein
    Obstruktions-Theorem: jede Fortsetzung der divergenten Prim-Spuren
    $\operatorname{Tr}T_s=P(s)$, $\operatorname{Tr}T_s^2=P(2s)$ bis $\Re s=\tfrac12$
    importiert $\zeta$ und seine Nullstellen via
    $P(s)=\sum_n \mu(n)/n\cdot\log\zeta(ns)$ --- also zirkulär. (Projekt-intern;
    kein literatur-verifizierter Standard.)
  \item \textbf{Toeplitz-Symbol-Submethode (Zweig~A).} $A_{\sin}$ ist nicht Toeplitz
    ($m$-abhängige oszillierende Diagonalen $\Rightarrow$ kein wohldefiniertes Symbol
    $f(\theta)$); ein struktureller Defekt des konstruierten Objekts, nicht bloß
    ,,die Methode gleicht dem offenen Problem``. Nachfolger: die endliche
    Sylvester-Inertia-Schranke (erledigt) plus GLT-/Outlier-Kontrolle.
  \item \textbf{Relative-Determinante-/Krein-Bypass.} Die relative Determinante
    $-\partial_z\log[\det(A-z)/\det(A^0-z)]$ entfernt den Euler-Pol-Kamm \emph{nicht}:
    ist $A^0$ der prolate Hintergrund, so ist er selbst kammfrei (nichts zu kürzen,
    und $\det(\text{Euler}\cdot\text{prolat})/\det(\text{prolat})=\det(\text{Euler})$
    behält den Kamm); die einzige kammfreie Alternative braucht einen
    selbstadjungierten Prim-Faktor, der zirkulär ist ($=$ RH). Dies ist eine ehrliche
    Selbstkorrektur einer intern vorgeschlagenen ,,richtigen Technologie``; der
    Vorgeschlagen-dann-widerlegt-Bogen ist in \cref{sec:sammler-disziplin}
    festgehalten.
\end{enumerate}

\noindent
Entscheidend: der \emph{Determinanten-Kanal} $-D'/D$ als \emph{Schärfung} der
CCM-Frontlinie bleibt gültig (\cref{sec:independent-results}); nur die
\emph{relative} Determinante als \emph{Bypass} der Fortsetzungs-Wand ist widerlegt.
Beide dürfen nicht verwechselt werden.


% =============================================================================
\section{Ruhende Routen}
\label{sec:dormant}
% =============================================================================

Eine kleine Anzahl konstruierter Routen ist ruhend statt tot: es handelt
sich um explizite Reduktionen oder Teilkonstruktionen, die nicht widerlegt
wurden, aber derzeit keine aktive Entwicklung erfahren. Das interne Register
(\texttt{AKTIONSPLAN.md}) führt das vollständige Inventar; die folgenden
vier werden hier zusammengefasst, da jede eine direkte Brücke zu einem
der aktiven Zweige aufweist.

\begin{itemize}
  \item \textbf{D-CCM (Dirichlet--CCM-Operator
    $D^{(\lambda,N,\chi)}_{\log}$).} Eine direkte $\chi$-Twistung der
    CCM-Rang-$N$-Störungskonstruktion. Adressiert GRH, nicht RH
    direkt --- der triviale Charakter $\chi_0$ (Riemann-Fall) ist in
    diesem Skalierungsregime praktisch nicht testbar. Begleiter von
    Zweig~Z.

  \item \textbf{D-Twist ($\chi$-getwistete Defektnorm).} Diagnostische
    Blaupause mit abgeschlossenem Pilot~v2 und Milestone~2$_\chi$-Arbeit;
    umformuliert als $\lambda$-asymptotische Aussage. Erste Tests bei
    $\chi_{21}, \chi_{33}$; Riemann in dieser Skala unerreichbar.

  \item \textbf{D-PW (Paley--Wiener-Sektordifferenz
    $\psi_\chi = \phi_\chi^+ - \phi_\chi^-$).} Klassifikation der
    fünf Paley--Wiener-Schritte abgeschlossen: S1, S2, S4 mechanisch;
    S3, S5 charakterspezifisch. Fortsetzung des veröffentlichten
    Dirichlet Character Atlas v2 (Zenodo~\texttt{20241612}).

  \item \textbf{P-M (Möbius-graduierter $\mathbb{Z}_2$-Primzahlhub).}
    Bedingtes Lemma für endliche Primzahlmenge bewiesen (Möbius-Identität
    exakt); volle Primzahlmenge konvergiert nur für $\Re(s)>1$;
    $\mathbb{Z}_2$-Sektoraussage für nichttriviale Nullstellen offen.
    Strukturell ansprechende Alternativinterpretation des Prime Orbit
    Axioms (Teil~I, Definition~2.1). Ein Obstruktions-Theorem 2026 schärft
    dies (\cref{sec:four-closures}): es existiert kein nicht-zirkuläres $L_s^-$,
    da die divergenten Prim-Spuren nur durch Import von $\zeta$ und seinen
    Nullstellen bis $\Re s=\tfrac12$ fortgesetzt werden können. Dies re-typisiert
    P-M als die \emph{Prim-Hälfte} einer P-M$+$M-Fusion --- ihr fehlt beweisbar die
    archimedisch-modulare Hälfte ($s\leftrightarrow 1-s$), die Route~M liefert
    (projekt-intern).
\end{itemize}

Die übrigen ruhenden Einträge (P-A, P-C, P-B, M, R, H) werden hier nicht
einzeln aufgeführt; siehe \texttt{AKTIONSPLAN.md} für das vollständige Inventar.
Zwei davon werden durch die Kritikalitäts-Linse aus \cref{sec:lambda-lens}
\emph{strategisch aufgewertet}: \textbf{M} (Bost--Connes\,$\leftrightarrow$\,Meyer)
und die globale adelische Trace-Route \textbf{P-A} (Connes--Consani-Scaling-Site)
sind vom exakt-strukturellen Typ und erhalten nun zusätzliches strategisches
Gewicht; \textbf{P-C} wird durch einen aktiven Zweig mit-gelöst; \textbf{P-B}
(Berry--Keating) ist vom richtigen Typ, aber unkonstruiert (Selbstadjungiertheit
ist das fehlende Hilbert--P\'olya-Theorem); \textbf{R}, \textbf{H} sind
Framework/expositorisch.


% =============================================================================
\section{Externe Linien}
\label{sec:external}
% =============================================================================

Die Landschaft liegt innerhalb eines weiteren Forschungsökosystems. Die
folgenden Linien werden als Ankerpunkte zitiert; nur die direkt
beweisrelevanten Einträge erscheinen hier.

\begin{itemize}
  \item \textbf{Connes 2026 Mainline} \cite{Connes2026}
    (arXiv:2602.04022). Reformuliert die Riemannschen Nullstellen als
    Infrarotspektrum eines selbstadjungierten Operators und identifiziert
    die zentrale offene Hypothese MS2 (\S6.6(b)). Die Zweige~A, B, Z
    konditionieren alle auf diese Linie.

  \item \textbf{Connes--van~Suijlekom 2025} \cite{ConnesvS2025}
    (arXiv:2511.23257). Theorem~5.6: der 1-dimensionale gerade Kern der
    trunkierten Weil-Form impliziert reelle Nullstellen von $\hat\eta$.
    Dies ist A1~=~B1 in der vorliegenden Landschaft.

  \item \textbf{Connes--Consani--Moscovici 2025}
    \cite{ConnesMoscovici2023} (arXiv:2511.22755). Grundlage der
    Rang-eins-Störung $D^{(\lambda,N)}_{\log}$. Verankert Zweig~B's
    Foundation-Lock und Zweig~Z's Microcluster-Konstruktion.

  \item \textbf{Groskin 2026} (arXiv:2605.20224). Unabhängige
    Implementierung der CCM-Galerkin-Matrix bei sechzehn Cutoffs
    $c=13\ldots 67, 100$ bei hoher Präzision (Aitken-Extrapolation,
    $\sim 10^{-334}$ Tiefe). Dient als externer Validator für die
    Zweige~B und~Z.

  \item \textbf{\'Sliwi\'nski 2026} (arXiv:2601.12133). Gegenregime:
    in der $\lambda = N$-Konfiguration mit niedriger Präzision wird eine
    invers-logarithmische Dissonanz-Unterschranke bewiesen,
    $\epsilon(\lambda,N)\ge 1/(4\ln\lambda)$. Schränkt das
    Fixed-$\lambda$-, $N\gg\lambda$-, Hochpräzisions-Setup dieser
    Serie nicht ein; steht als Warnung gegen naive $\lambda = N$-Limites.

  \item \textbf{Watson 2026} (arXiv:2602.01248), ,,The Riemann $\Xi$-function from
    primitive Markovian cycles~I.`` Konstruiert eine Laguerre--P\'olya-Funktion
    $\Psi$ aus endlicher reversibler Markov-Dynamik, ,,entirely Archimedean, using
    no Euler products, primes, or arithmetic spectral input``; kein RH-Beweis --- es
    isoliert ,,a single remaining analytic problem relating $\Psi$ to
    $\Xi(2\cdot)$``. Eine unabhängige Triangulation der gemeinsamen Wand: von der
    archimedischen Seite erreicht es denselben \emph{Identifikations}-Schritt (das
    konstruierte Objekt mit $\Xi$ zu identifizieren), den die Routen dieser Serie
    von der arithmetischen Seite erreichen (\cref{sec:common-wall-ac}).

  \item \textbf{Franchi Vicer\'e 2026} (Zenodo~19546495). Beansprucht
    einen unbedingten RH-Beweis über Apeiroeder / semi-lokale spektrale
    Abstiegsmethode. Untersucht und als nicht tragfähig befunden: das
    Stabilisierungsargument verwechselt punktweise mit
    streifenuniformer Konvergenz; der Apeiroeder-Limes-Schritt ist nicht
    rigoros; Randbedingungen ersetzen nicht den Even-Sector / MS1-Input.
    Siehe \texttt{AKTIONSPLAN.md} \S Klasse~8.
\end{itemize}


% =============================================================================
\section{Systematisch erkundete Alternativen}
\label{sec:explored-alternatives}
% =============================================================================

Elf alternative Strategien (BI-1 bis BI-11) wurden während der
Entwicklung systematisch erkundet. Diese erschöpfende Suche ist selbst ein
Meta-Ergebnis: der Lösungsraum ist kartiert.

\begin{center}
\footnotesize
\renewcommand{\arraystretch}{1.2}
\resizebox{\textwidth}{!}{%
\begin{tabular}{l|p{4.8cm}|l|l}
\textbf{Ansatz} & \textbf{Kernidee} & \textbf{Ergebnis} & \textbf{Status} \\\hline
BI-1: $\det_2$ / Carleman & Auf Hilbert--Schmidt-Klasse erweitern & Kein Positivitätsersatz & Widerlegt \\
BI-2: RG-Fixpunkt & Weil-Positivität unter Skalierung & Nicht formalisierbar & Verworfen \\
BI-3: Universelle Eigenschaft & Physikalische Plausibilität $\Rightarrow$ RH & Zu abstrakt & Verworfen \\
BI-5: Liftprogramme & Weil $\to$ Grothendieck $\to$ RMT & Alle offen außer $\mathbb{F}_q$ & Nicht nutzbar \\
BI-6: Implikationsring & Li $\leftrightarrow$ Weil $\leftrightarrow$ NB $\leftrightarrow$ BD & Ring existiert nicht & Widerlegt \\
BI-7: Pöschl--Teller-Streuung & Gamma-Struktur analog zu $\xi$ & Keine Off-Line-Stabilität & Negativ \\
\textbf{BI-8: Stabilität zweiter Ordnung} & \textbf{RH als Variationsproblem} & \textbf{Architektur korrekt} & \textbf{Realisiert} \\
BI-9: Solitonenmodelle & Topologische Stabilität & Coulomb-Gas instabil & Widerlegt \\
\textbf{BI-10: Bausteine} & \textbf{No-Coordination, Shift-Struktur} & \textbf{Erfolgreich wiederverwendet} & \textbf{Integriert} \\
\textbf{BI-11: Connes 2026} & \textbf{Even Dominance über $Q_W$} & \textbf{Durchbruch} & \textbf{Kernpfad}
\end{tabular}}
\end{center}

\noindent Der gewählte Pfad (BI-8 + BI-10 + BI-11) kombiniert die
Variationsarchitektur, die arithmetischen Bausteine und Connes'
spektrales Framework. Die Routen~A und~B (totale Positivität,
kommutierender Operator) waren Ausschlüsse, die strukturelle
Entdeckungen darstellen: Even Dominance gilt \emph{trotz} ihrer
Abwesenheit, was verdeutlicht, dass der Mechanismus genuinen
arithmetisch-statistischen und nicht algebraischen Charakter hat
(Teil~II~\cite{Geiger2026partII}).


% =============================================================================
\section{Unabhängige Ergebnisse}
\label{sec:independent-results}
% =============================================================================

Mehrere Ergebnisse dieses Programms sind von unabhängigem Interesse,
ungeachtet des Status des Hauptbeweises:

\begin{center}
\footnotesize
\renewcommand{\arraystretch}{1.2}
\resizebox{\textwidth}{!}{%
\begin{tabular}{p{4.2cm}|p{3.2cm}|p{4.2cm}|l}
\textbf{Ergebnis} & \textbf{Kontext} & \textbf{Bedeutung} & \textbf{Status} \\\hline
\multicolumn{4}{l}{\emph{Positive Ergebnisse}} \\[2pt]
Shift Parity Lemma & Arithmetische Trigonometrie & Jede Primzahl bevorzugt gerade Eigenfunktionen & Bewiesen \\
No-Coordination Lemma & Multiplikative Unabhängigkeit & Arithmetischer Ersatz für Entropieterm & Bewiesen \\
Leading-Mode Cancellation ($c = 2+\sqrt{2}$) & Resolventenanalyse & Universelle paarweise Auslöschungskonstante & Bewiesen \\
Euler--Maclaurin: $\rho^{\mathrm{EM}} < 0$ & PNT-Asymptotik & Globaler Stabilitätsanker & Bewiesen \\
$33$ endliche CAP-Diagnostiken & Intervallarithmetik & Bedingte Brücke & Evidenz \\
Konzentrations-Lemma & Even-Dominance-Bildmaß & Keine feste Primzahlmenge kontrolliert einen positiven Anteil der Primmasse & Bewiesen \\
Endliche Sylvester-Inertia ($\nu_-=0$) & Intervallarithmetik, $\lambda=100$ & Finite Even Dominance rigoros, kein Symbol nötig & Verifiziert \\
Residue-Calibration-Kanal & Weyl-Funktions-Diagnostik & Eigenwert-Konv.\ $\neq$ Weyl-Konv.; Kanal ist $-D'/D$, $\det_{\mathrm{reg}}\!\to C\Xi$ & Diagnostik \\[4pt]
\multicolumn{4}{l}{\emph{Negative Ergebnisse}} \\[2pt]
Obstruktionen K1--K4 & Sackgassen-Analyse & Kartiert unmögliche Strategien & Dokumentiert \\
Kein absolutes $\mathrm{PF}_\infty$ für $A_\lambda$ & Primzahl-Shift-Multiplikator & $M(\xi)$ negativ auf $\sim49\,\%$ des Definitionsbereichs & Bewiesen \\
Kein universeller kommutierender Operator & Sturm--Liouville-Suche & Kern von $[T, D_N(r)]$ ist $\{c \cdot I\}$ & Bewiesen \\[4pt]
\multicolumn{4}{l}{\emph{Strukturelle Einsichten}} \\[2pt]
Invertierte-Hesse-Struktur & YM/TU/RH-Vergleich & RH ist ein Maximum, kein Minimum & Bestätigt \\
Interdisziplinärer Katalog & Regularisierung & $11$ Felder teilen Spurklassen-Obstruktion & Dokumentiert
\end{tabular}}
\end{center}

\noindent Die negativen Ergebnisse schützen zukünftige Forschende davor,
Sackgassen-Strategien zu wiederholen. Das Shift Parity Lemma, die
Obstruktionsanalyse (K1--K4) und die Mellin-Abklingbarriere sind
unabhängig publizierbar.


% =============================================================================
\section{Hypothesen-Korrekturen im iterativen Forschungsprozess}
\label{sec:sammler-disziplin}
% =============================================================================

Die iterative Entwicklung brachte mehrere Hypothesen-Korrekturen hervor,
die hier als Dokumentation ehrlicher Forschungspraxis festgehalten werden
(Sammler-Disziplin).

\subsection{GF1-Korrekturkette (Iterationen 15--29)}

\begin{center}
\renewcommand{\arraystretch}{1.2}
\resizebox{\textwidth}{!}{%
\begin{tabular}{c|l|l}
\toprule
\textbf{Iter.} & \textbf{Hypothese} & \textbf{Status} \\
\midrule
15 & $C_m \approx 8$ universell & falsifiziert (Randeffekt) \\
18 & $C_m(\lambda=5) \approx 4{,}3$ & falsifiziert (kein Plateau) \\
19 & $\mathrm{sLI} \sim 1/N$ & bestätigt, aber ohne $\lambda$-Skalierung \\
20 & $N \cdot \mathrm{sLI} \approx 329$ universell & falsifiziert \\
\textbf{27} & $\mathbf{\mathrm{sLI} = C^*\,\lambda/(NL) + o(\cdot)}$ & bestätigt für $\lambda=5,7$, $C^* = 64\pi^2/3$ \\
28 & $\lambda=3$ All-$\lambda$-Prüfung & Niedrigband-Ausreißer; Zwei-Regime-Guardrail \\
29 & $\lambda=11$-Diskriminator und Endpoint-Kern-Audit & $C_{\mathrm{eff}}\approx168$ bei $N=200$; einfache Endpoint-Faktorisierung falsifiziert \\
\bottomrule
\end{tabular}
}
\end{center}


\subsection{\texorpdfstring{B-1-Korrekturen und RH-neutrale Form von $c_\Xi$}{B-1-Korrekturen und RH-neutrale Form von cXi}}

Die Hadamard-Konstante $c_\Xi$ wurde zunächst in der RH-spezialisierten
Form $\sum_{k\geq 1} 1/\gamma_k^2$ geschrieben (Summe über die
Imaginärteile der Nullstellen auf der kritischen Geraden). Ein externes
Audit korrigierte dies zur RH-neutralen Form
$c_\Xi = -\tfrac{1}{2}\sum_\rho 1/(\rho-\tfrac{1}{2})^2$
(Summe über alle nichttrivialen Nullstellen). Dies ist eine
Sammler-Disziplin-Korrektur: die RH-neutrale Form ist mathematisch
allgemeiner und setzt nicht voraus, was bewiesen werden soll.


\subsection{Der Relative-Determinante-Bypass: vorgeschlagen und widerlegt (2026)}

In der Phase 2026 wurde die relative Determinante (Krein-Spektral-Shift)
$-\partial_z\log[\det(A-z)/\det(A^0-z)]$ intern als die ,,richtige Technologie``
vorgeschlagen, um den Euler-Pol-Kamm in der oberen Halbebene zu entfernen und so die
Fortsetzungs-Wand zu umgehen (\cref{rem:wall-folklore}). Ein anschließender Test am
echten Connes--Moscovici-Prolaten-Hintergrund widerlegte dies als Bypass: der
prolate $A^0$ ist selbst Herglotz (kammfrei), sodass nichts zu kürzen ist
($\det(\text{Euler}\cdot\text{prolat})/\det(\text{prolat})=\det(\text{Euler})$ behält
den Kamm), und die einzige kammfreie Alternative braucht einen selbstadjungierten
Prim-Faktor, der zirkulär ist ($=$~RH). Der Vorschlag wird daher als Bypass
zurückgezogen. Der \emph{separate} Befund --- dass der Determinanten-Kanal $-D'/D$,
nicht bloße Eigenwert-Konvergenz, das korrekte zu konvergierende Objekt ist
($\det_{\mathrm{reg}}(A-z)\to C\,\Xi$) --- bleibt gültig
(\cref{sec:independent-results}); beide dürfen nicht vermengt werden.

Eine zweite Korrektur: ein internes Phasen-Lemma
$\Theta_\sigma=\xi(\sigma-\mathrm{i}z)/\xi(\sigma+\mathrm{i}z)$ war vorzeichenfalsch
(seine Pole kommen aus dem Nenner, sodass es unter RH Pole \emph{hat}); das korrekte
Objekt ist die reziproke Phase
$\widehat\Theta_\sigma=\xi(\sigma+\mathrm{i}z)/\xi(\sigma-\mathrm{i}z)$, mit
RH~$\Leftrightarrow\widehat\Theta_\sigma$ polfrei in der oberen Halbebene für alle
$\sigma>\tfrac12$.


\subsection{Wert der Korrektur-Dokumentation}

Diese Korrekturen zeigen, dass der Forschungsprozess selbstkorrigierend
ist. Jede Korrektur schärfte das theoretische Gerüst: die GF1-Kette
bewegte sich von einer naiven Konstante ($C_m \approx 8$) über mehrere
Zwischenformen zum $\lambda=5,7$-Boundary-Moment-Benchmark
$\mathrm{sLI} \sim C^*\lambda/(NL)$; der $\lambda=3$-Sweep und der offene
$\lambda=11$-Diskriminator verhindern eine All-$\lambda$-Satzlesart, und die $c_\Xi$-Korrektur entfernte eine
zirkuläre RH-Abhängigkeit aus einer Schlüsselkonstante.


% =============================================================================
\section{Zusammenfassung}
\label{sec:summary}
% =============================================================================

Teil~IV liefert die Cross-Route-Perspektive: die drei aktiven Zweige
teilen ein Galerkin-Muster (Fixed-Parameter positiv, Uniform-Limit offen),
und zwei isolierte offene Beweisslots (B-1 Log-Krümmung, GF1
Faktor-16-Identifikation) sind die konkreten analytischen Ziele, die
wohlgestellte offene Probleme bleiben. \Cref{sec:sharpened-wall} ergänzt dann eine
\emph{heuristische} Kritikalitäts-Linse ($\mathrm{RH}\Leftrightarrow\Lambda=0$), die
die Routen zu exakt-strukturellen Mechanismen hin neu gewichtet (ohne die
Marge-artigen Slots zu widerlegen), lokalisiert die gemeinsame Wand als die
analytische Fortsetzung über $\Re s=1$, dokumentiert eine berechenbare
Laguerre--P\'olya/Herglotz-Diagnostik (ein klassisches Kriterium als
Test-Instrument) und berichtet vier Schließungen \emph{konkreter Konstruktionen}
--- die OS-Reflexion in natürlicher Form, das naive Möbius-$L_s^-$, die
Toeplitz-Symbol-Submethode und den relative-Determinante-Bypass --- keine davon
schließt eine Route oder RH. Die
ruhenden und externen Routen werden der Vollständigkeit halber
katalogisiert; die Hypothesen-Korrektur-Dokumentation liefert Transparenz.
Der komprimierte Atlas und die Gesamtbewertung folgen in
Teil~V~\cite{Geiger2026partV}.


\bigskip
\noindent\textbf{Danksagungen.} Berechnungen wurden auf einem
Mac~Studio (Apple M1~Max, 32~GB) und einem Hetzner CCX13 Cloud-Server
(2~dedizierte vCPU, 8~GB) durchgeführt. Die Intervallarithmetik-Bibliothek
mpmath wurde von Fredrik Johansson entwickelt.


\section*{KI-Offenlegung}

Dieses Manuskript wurde in umfangreicher Zusammenarbeit mit
KI-Sprachmodellen entwickelt (Claude, Anthropic; Copilot,
Microsoft/GitHub; Gemini, Google DeepMind). KI trug zur konzeptuellen
Exploration, analytischen Arbeit, Literaturrecherche, zum Schreiben und
zur kritischen Überprüfung bei. Der Autor, Lukas Geiger, konzipierte die
Forschungsrichtung, brachte Domänenexpertise ein und übte die endgültige
redaktionelle Autorität aus. Der Autor übernimmt die volle und alleinige
wissenschaftliche Verantwortung für alle Inhalte, einschließlich
etwaiger Fehler.


\bibliographystyle{plain}
\bibliography{fst-rh-references}

\end{document}
